\begin{abstract}
    Dealing with long and highly complex technical text is a challenge for Large Language Models (LLMs), which still have to unfold their potential in supporting expensive and time-intensive processes like patent drafting.
    Within patents, the description constitutes more than 90\% of the document on average.
    Yet, its automatic generation remains understudied.
    When drafting patent applications, patent attorneys typically receive invention reports (IRs), which are usually confidential, hindering research on LLM-supported patent drafting.
    Often, pre-publication research papers serve as IRs.
    We leverage this duality to build \DatasetName, an open and realistic benchmark for patent drafting consisting of 1.8k patent-paper pairs describing the same inventions.
    To address the complex long-document patent generation task, we propose chunk-based outline-guided generation using the research paper as technical specification of the invention.
    Our extensive evaluation using \DatasetName and a human case study show that LLMs can effectively leverage information from the paper, but still struggle to provide the necessary level of detail. 
    Fine-tuning leads to more patent-style language, but also to more hallucination. 
    We release our data and code at \url{https://github.com/boschresearch/Pap2Pat}.
\end{abstract}